%-----------------------------------------------------------------------
% Reporte LogGuard V2 — Autómata Finito No Determinista para
% Detección de Amenazas en Logs de PostgreSQL
% Autor: Adolfo Toledo
%-----------------------------------------------------------------------
\documentclass[12pt, a4paper]{report}

%---------- Paquetes ---------------------------------------------------
\usepackage[utf8]{inputenc}
\usepackage[T1]{fontenc}
\usepackage[spanish, es-noquoting]{babel}
\usepackage{geometry}
\usepackage{graphicx}
\usepackage{amsmath, amssymb, amsthm}
\usepackage{booktabs}
\usepackage{longtable}
\usepackage{array}
\usepackage{xcolor}
\usepackage{listings}
\usepackage{tikz}
\usepackage{fancyhdr}
\usepackage{titlesec}
\usepackage{caption}
\usepackage{subcaption}
\usepackage{multirow}
\usepackage{tabularx}
\usepackage{enumitem}
\usepackage{float}
\usepackage[hypertexnames=false]{hyperref}

\usetikzlibrary{automata, positioning, arrows.meta, fit, decorations.pathreplacing}

\geometry{
  top=2.5cm, bottom=2.5cm,
  left=3cm, right=2.5cm
}
\setlength{\headheight}{15pt}
\setlength{\emergencystretch}{3em}
\hbadness=10000
\vbadness=10000

%---------- Colores ----------------------------------------------------
\definecolor{codeblue}{RGB}{0,80,180}
\definecolor{codegray}{RGB}{128,128,128}
\definecolor{codegreen}{RGB}{0,140,0}
\definecolor{codebg}{RGB}{248,248,250}
\definecolor{accent}{RGB}{30,90,200}
\definecolor{danger}{RGB}{200,30,30}
\definecolor{success}{RGB}{30,160,60}

%---------- Listings ---------------------------------------------------
\lstdefinestyle{json}{
  basicstyle=\ttfamily\footnotesize,
  backgroundcolor=\color{codebg},
  breaklines=true,
  frame=single,
  rulecolor=\color{codegray!50},
  numbers=left,
  numberstyle=\tiny\color{codegray},
  keywordstyle=\color{codeblue}\bfseries,
  stringstyle=\color{codegreen},
  commentstyle=\color{codegray}\itshape,
  tabsize=2,
  showstringspaces=false,
  captionpos=b
}

\lstdefinestyle{csharp}{
  language={[Sharp]C},
  basicstyle=\ttfamily\footnotesize,
  backgroundcolor=\color{codebg},
  breaklines=true,
  frame=single,
  rulecolor=\color{codegray!50},
  numbers=left,
  numberstyle=\tiny\color{codegray},
  keywordstyle=\color{codeblue}\bfseries,
  stringstyle=\color{codegreen},
  commentstyle=\color{codegray}\itshape,
  tabsize=2,
  showstringspaces=false,
  captionpos=b
}

%---------- Encabezados ------------------------------------------------
\pagestyle{fancy}
\fancyhf{}
\fancyhead[L]{\small LogGuard V2 — Reporte Técnico}
\fancyhead[R]{\small Adolfo Toledo}
\fancyfoot[C]{\thepage}
\renewcommand{\headrulewidth}{0.4pt}

%---------- Teoremas ---------------------------------------------------
\theoremstyle{definition}
\newtheorem{definicion}{Definición}[chapter]
\newtheorem{ejemplo}{Ejemplo}[chapter]
\theoremstyle{remark}
\newtheorem{nota}{Nota}[chapter]

%---------- Comandos auxiliares ----------------------------------------
\newcommand{\tok}[1]{\texttt{\color{codeblue}#1}}
\newcommand{\state}[1]{\ensuremath{q_{\text{#1}}}}
\newcommand{\nfa}[1]{\ensuremath{\mathcal{#1}}}
\newcommand{\severity}[2]{\textbf{\color{#1}#2}}

%=======================================================================
\begin{document}
\pagenumbering{roman}
%=======================================================================

%----------------------------------------------------------------------
% PORTADA
%----------------------------------------------------------------------
\begin{titlepage}
  \centering

  \vspace*{1.5cm}

  {\Huge\bfseries\color{accent} LogGuard V2}\\[0.4cm]
  {\Large Reporte Técnico de Autómatas Finitos No Deterministas}\\[0.3cm]
  {\large Detección de Amenazas en Logs de PostgreSQL}

  \vspace{1cm}
  \rule{12cm}{0.8pt}
  \vspace{1cm}

  \begin{minipage}{0.85\textwidth}
    \centering
    {\normalsize
    Este documento presenta la descripción formal del lenguaje de patrones de
    ataque reconocido por LogGuard V2, incluyendo la definición del sistema
    de autómatas finitos no deterministas, diagramas de estados, tablas de
    transiciones, decisiones de diseño y análisis de cadenas de prueba.
    }
  \end{minipage}

  \vspace{2.5cm}

  \begin{tabular}{rl}
    \textbf{Autor:}    & Adolfo Toledo \\[0.3cm]
    \textbf{Versión:}  & 2.1.0 \\[0.3cm]
    \textbf{Fecha:}    & Junio 2026 \\[0.3cm]
    \textbf{Plataforma:} & .NET 10.0 / Windows / WPF \\[0.3cm]
    \textbf{Tecnología:} & Autómata Finito No Determinista (NFA) \\
  \end{tabular}

  \vspace{3cm}
  \rule{12cm}{0.4pt}
  \vspace{0.5cm}

  {\small Universidad --- Teoría de la Computación}

  \vfill
\end{titlepage}

%----------------------------------------------------------------------
% ÍNDICE
%----------------------------------------------------------------------
\tableofcontents
\clearpage
\pagenumbering{arabic}

%======================================================================
\chapter{Descripción del Lenguaje}
%======================================================================

\section{Contexto del Sistema}

LogGuard V2 es una aplicación de seguridad en tiempo real que monitorea archivos
de registro (\emph{logs}) de instancias PostgreSQL. Su objetivo es detectar
comportamientos maliciosos en consultas SQL antes de que causen daño. El sistema
lee logs en formato estándar de PostgreSQL, extrae las sentencias SQL, las
tokeniza y aplica un conjunto de Autómatas Finitos No Deterministas (NFA) para
clasificar cada consulta como \emph{benigna} o \emph{amenaza}.

\section{El Lenguaje de Entrada}

El sistema opera sobre dos capas de lenguaje:

\subsection{Capa 1 — Formato de Log de PostgreSQL}

Las líneas de log siguen una de dos plantillas de prefijo configurables en
\texttt{postgresql.conf} mediante el parámetro \texttt{log\_line\_prefix}:

\begin{center}
\small
\resizebox{\textwidth}{!}{%
\begin{tabular}{lll}
  \toprule
  \textbf{Estilo} & \textbf{Prefijo} & \textbf{Ejemplo} \\
  \midrule
  \texttt{\%m} & \texttt{\%m [\%p] \%q\%u@\%h \%d} &
    \texttt{2026-06-21 20:10:15.123 UTC [1234] app@10.0.0.1 mydb} \\
  \texttt{\%t} & \texttt{\%t [\%p]: [\%l-\%s] user=\%u,...} &
    \texttt{2026-06-21 20:10:15 UTC [1234]: [1-1] user=app,db=mydb} \\
  \bottomrule
\end{tabular}%
}
\end{center}

Cada línea contiene: \emph{timestamp} (ISO~8601 + zona horaria), PID del proceso,
nivel de severidad (\texttt{LOG}, \texttt{ERROR}, \texttt{FATAL}, etc.),
y el cuerpo del mensaje. Los tipos de línea reconocidos son:

\begin{description}[leftmargin=3cm, labelindent=0.5cm]
  \item[\texttt{Statement}] Consulta SQL ejecutada (\texttt{statement: <sql>}).
  \item[\texttt{Duration}]  Tiempo de ejecución en ms (\texttt{duration: X.Y ms}).
  \item[\texttt{ConnectionReceived}] Nueva conexión entrante con host y puerto.
  \item[\texttt{ConnectionAuthenticated}] Método de autenticación usado.
  \item[\texttt{ConnectionAuthorized}]   Sesión autorizada: usuario y base de datos.
  \item[\texttt{Disconnection}] Cierre de sesión con resumen de duración.
  \item[\texttt{General}] Mensaje de log sin campos estructurados adicionales.
  \item[\texttt{Unknown}] Línea no reconocida por ningún patrón.
\end{description}

\subsection{Capa 2 — Lenguaje de Tokens SQL}

Una vez extraída la sentencia SQL del campo \texttt{statement:}, el sistema la
transforma mediante un proceso de cuatro fases en una \emph{secuencia de tokens}.
Esta secuencia es el \textbf{lenguaje de entrada formal} para los autómatas.

El alfabeto $\Sigma$ del lenguaje está compuesto por más de 70 categorías de
tokens, agrupadas en:

\begin{center}
\begin{tabularx}{\textwidth}{lX}
  \toprule
  \textbf{Categoría} & \textbf{Tokens representativos} \\
  \midrule
  Estructura de consulta & \tok{SELECT}, \tok{FROM}, \tok{WHERE}, \tok{JOIN},
                           \tok{UNION}, \tok{UNION\_ALL}, \tok{GROUP}, \tok{ORDER} \\
  Comparación lógica     & \tok{OR}, \tok{AND}, \tok{NOT}, \tok{IN}, \tok{LIKE},
                           \tok{EQUALS}, \tok{NEQ}, \tok{GT}, \tok{LT} \\
  DML                    & \tok{INSERT}, \tok{UPDATE}, \tok{DELETE}, \tok{MERGE},
                           \tok{VALUES}, \tok{RETURNING} \\
  DDL / Privilegios      & \tok{CREATE}, \tok{DROP}, \tok{ALTER}, \tok{GRANT},
                           \tok{REVOKE}, \tok{SUPERUSER} \\
  Funciones peligrosas   & \tok{SLEEP}, \tok{BENCHMARK}, \tok{WAITFOR\_DELAY},
                           \tok{XP\_CMDSHELL}, \tok{LOAD\_FILE} \\
  Tablas del sistema     & \tok{INFORMATION\_SCHEMA}, \tok{PG\_SHADOW},
                           \tok{PG\_USER}, \tok{PG\_ROLES} \\
  Especiales             & \tok{TAUTOLOGY}, \tok{COMMENT}, \tok{STAR} \\
  Primitivos             & \tok{STRING}, \tok{NUMBER}, \tok{IDENT},
                           \tok{SEMICOLON}, \tok{LPAREN}, \tok{RPAREN} \\
  \bottomrule
\end{tabularx}
\end{center}

El token \tok{TAUTOLOGY} no proviene del léxico SQL estándar sino que es
generado en la \emph{Fase~2} del tokenizador cuando se detecta una expresión
lógicamente siempre verdadera (e.g., \texttt{1=1}, \texttt{'x'='x'}, \texttt{a=a}).

\section{Pipeline de Procesamiento}

El flujo completo desde el archivo de log hasta la alerta es:

\begin{enumerate}
  \item \textbf{Lectura reactiva:} \texttt{FileWatcherLive} lee el archivo con
        sondeo cada 500~ms, detecta rotación por cambio de \emph{inode} y
        maneja líneas de continuación multi-línea.
  \item \textbf{Parseo de log:} \texttt{PostgreSqlLogParser} aplica 7 expresiones
        regulares precompiladas para clasificar cada línea y extraer metadatos
        estructurados.
  \item \textbf{Tokenización SQL:} \texttt{SqlTokenizer} ejecuta 4 fases secuenciales.
  \item \textbf{Evaluación NFA:} \texttt{NfaEngine} ejecuta todos los perfiles NFA
        habilitados sobre la secuencia de tokens.
  \item \textbf{Alerta:} Si algún NFA acepta, se genera una entrada en la bitácora
        de amenazas con nivel de severidad y tipo de ataque.
\end{enumerate}

\section{Fases del Tokenizador}

\begin{description}[leftmargin=3.5cm]
  \item[Fase 1 — Normalización] Se eliminan comentarios SQL (\texttt{--},
        \texttt{/* */}), se decodifican literales hexadecimales
        (\texttt{0x41} $\to$ \texttt{A}) y se unifican espacios blancos.
  \item[Fase 2 — Detección de tautologías] Se aplican 5 expresiones regulares
        acotadas para identificar expresiones siempre verdaderas y sustituirlas
        por el marcador \tok{TAUTOLOGY}.
  \item[Fase 3 — Escaneo léxico] Se recorre el texto normalizado carácter a
        carácter y se emiten tokens usando un diccionario de 160 entradas y
        expresiones regulares para literales.
  \item[Fase 4 — Fusión] Pares de tokens adyacentes se combinan donde
        corresponde (e.g., \tok{UNION} + \tok{ALL} $\to$ \tok{UNION\_ALL}).
\end{description}

%======================================================================
\chapter{Definición Formal}
%======================================================================

\section{Autómata Finito No Determinista (NFA)}

\begin{definicion}[Autómata Finito No Determinista]
Un NFA se define como una 5-tupla
\[
  \nfa{M} = (Q,\; \Sigma,\; \delta,\; q_0,\; F)
\]
donde:
\begin{itemize}
  \item $Q$ es el conjunto finito de estados,
  \item $\Sigma$ es el alfabeto finito de símbolos de entrada (tokens SQL),
  \item $\delta : Q \times \Sigma \to \mathcal{P}(Q)$ es la función de
        transición no determinista,
  \item $q_0 \in Q$ es el estado inicial,
  \item $F \subseteq Q$ es el conjunto de estados de aceptación.
\end{itemize}
\end{definicion}

\begin{definicion}[Lenguaje aceptado]
El lenguaje aceptado por $\nfa{M}$ es
\[
  L(\nfa{M}) = \{ w \in \Sigma^* \mid \hat\delta(\{q_0\}, w) \cap F \neq \emptyset \}
\]
donde $\hat\delta$ es la extensión de $\delta$ a cadenas.
\end{definicion}

\begin{definicion}[Subcadena aceptada — modo substring]
LogGuard opera en modo de reconocimiento de \emph{subcadena}: el estado inicial
$q_0$ se reinyecta al conjunto activo en cada paso, de modo que el autómata
detecta el patrón en cualquier posición dentro de la secuencia de tokens:
\[
  \text{match}(w) = \exists\, 0 \le i \le j \le |w| : \hat\delta(\{q_0\}, w[i..j]) \cap F \neq \emptyset
\]
\end{definicion}

\section{Perfil de Amenaza}

Cada perfil NFA en LogGuard amplía la definición estándar con dos elementos:

\begin{definicion}[Perfil NFA con restricción negativa]
Un perfil NFA extendido es una 7-tupla
\[
  \nfa{P} = (Q,\; \Sigma,\; \delta,\; q_0,\; F,\; T,\; \mathit{absent})
\]
donde:
\begin{itemize}
  \item $T$ es el tipo de amenaza (\textsc{sqli}, \textsc{exfil},
        \textsc{privesc}, \textsc{bruteforce}, \textsc{discovery}),
  \item $\mathit{absent} \subseteq \Sigma$ es el conjunto de tokens cuya
        \emph{presencia} en $w$ invalida la aceptación.
\end{itemize}
La aceptación final se define:
\[
  \text{accept}(w) = \bigl[\hat\delta(\{q_0\},w) \cap F \neq \emptyset\bigr]
                   \;\wedge\; \forall\, a \in \mathit{absent} : a \notin w
\]
\end{definicion}

\section{Algoritmo de Simulación}

El algoritmo implementado en \texttt{NfaEngine.Run} es una simulación de
subconjuntos incremental de complejidad $O(|w| \cdot |Q|)$:

\begin{lstlisting}[style=csharp, caption={Núcleo del motor NFA (NfaEngine.cs)}]
HashSet<string> active = new() { startState };

foreach (string token in tokens)
{
    var next = new HashSet<string> { startState };  // modo substring

    foreach (string state in active)
        if (_delta.TryGetValue((state, token), out var targets))
            next.UnionWith(targets);

    active = next;

    if (active.Overlaps(_acceptStates))
        if (!tokens.Intersect(profile.RequireAbsentTokens).Any())
            return true;
}
return false;
\end{lstlisting}

\section{Inventario de Perfiles NFA}

LogGuard incluye 6 perfiles NFA cargados desde archivos JSON en tiempo de
ejecución mediante \texttt{NfaLoader}:

\begin{center}
\small
\resizebox{\textwidth}{!}{%
\begin{tabular}{llccl}
  \toprule
  \textbf{Perfil} & \textbf{ID} & \textbf{Estados} & \textbf{Transiciones}
                  & \textbf{Severidad} \\
  \midrule
  SQL Injection      & \texttt{pgsql-sqli-v2}        & 10 & 29
    & \severity{danger}{High} \\
  Data Exfiltration  & \texttt{pgsql-exfil-v1}       & 5  & 6
    & \severity{danger}{High} \\
  Privilege Escal.   & \texttt{pgsql-privesc-v1}     & 6  & 7
    & \severity{danger}{Critical} \\
  Discovery/Enum.    & \texttt{pgsql-discovery-v2}   & 4  & 5
    & \severity{codegray}{Medium} \\
  Brute Force        & \texttt{pgsql-bruteforce-v1}  & 8  & 8
    & \severity{codegray}{Medium} \\
  Time-Based SQLi    & \texttt{pgsql-time-sqli-v2}   & 5  & 6
    & \severity{danger}{High} \\
  \bottomrule
\end{tabular}%
}
\end{center}

%======================================================================
\chapter{Diagramas de Autómatas}
%======================================================================

Esta sección presenta los diagramas de transición de estados para los tres
perfiles más relevantes: SQL Injection, Data Exfiltration y Privilege Escalation.

%----------------------------------------------------------------------
\section{NFA — Inyección SQL (\texttt{pgsql-sqli-v2})}

El autómata de inyección SQL reconoce cinco patrones de ataque:
tautología directa, bypass OR, UNION injection, funciones de tiempo
(\texttt{SLEEP}/\texttt{BENCHMARK}) y enumeración de esquemas
(\texttt{INFORMATION\_SCHEMA}).

\begin{figure}[htbp]
\centering
\begin{tikzpicture}[
  ->, >=Stealth,
  node distance=2.8cm,
  every state/.style={draw, circle, minimum size=1.0cm, font=\scriptsize},
  accepting/.style={double},
  initial text=,
  auto
]

% Estados
\node[state, initial]                    (q0)      {$q_0$};
\node[state, right of=q0]               (qpre)    {$q_p$};
\node[state, right of=qpre]             (qtbl)    {$q_t$};
\node[state, right of=qtbl]             (qwhere)  {$q_w$};
\node[state, below of=qwhere]           (qcol)    {$q_c$};
\node[state, left of=qcol]              (qeq)     {$q_e$};
\node[state, left of=qeq]               (qval)    {$q_v$};
\node[state, below of=q0, yshift=-1cm]  (qor)     {$q_{or}$};
\node[state, right of=qor, xshift=1.5cm](qunion)  {$q_u$};
\node[state, accepting,
      below of=qval, yshift=-0.5cm]     (qsqli)   {\textbf{$q_!$}};

% --- Transiciones desde q0 ---
\draw (q0)     edge[bend left=10]  node[above]{\scriptsize SELECT}     (qpre);
\draw (q0)     edge[bend left=40]  node[above, sloped]{\scriptsize TAUTOLOGY / SLEEP / IS} (qsqli);
\draw (q0)     edge               node[left]{\scriptsize OR}           (qor);
\draw (q0)     edge[bend right=10] node[below]{\scriptsize UNION/U\_A} (qunion);

% --- Camino principal ---
\draw (qpre)   edge[loop above]    node{\scriptsize IDENT/STAR}        (qpre);
\draw (qpre)   edge               node[above]{\scriptsize FROM}        (qtbl);
\draw (qpre)   edge[bend right=45] node[right, sloped]{\scriptsize SLEEP} (qsqli);

\draw (qtbl)   edge[loop above]    node{\scriptsize IDENT}             (qtbl);
\draw (qtbl)   edge               node[right]{\scriptsize WHERE}       (qwhere);
\draw (qtbl)   edge[bend right=30] node[above, sloped]{\scriptsize IS/SEMI} (qsqli);

\draw (qwhere) edge               node[right]{\scriptsize IDENT}       (qcol);
\draw (qwhere) edge[bend left=35] node[above, sloped]{\scriptsize TAUTO/SEMI} (qsqli);

\draw (qcol)   edge               node[below]{\scriptsize EQUALS}      (qeq);

\draw (qeq)    edge               node[below]{\scriptsize STR/NUM}     (qval);

\draw (qval)   edge               node[left]{\scriptsize OR/UNION/\ldots} (qsqli);

% --- Bypass OR ---
\draw (qor)    edge[bend right=10] node[below]{\scriptsize TAUTO/NUM/STR} (qsqli);

% --- UNION path ---
\draw (qunion) edge[bend right=10] node[above]{\scriptsize SELECT}     (qsqli);

% --- q_sqli loops ---
\draw (qsqli)  edge[loop right]    node{\scriptsize SELECT/TAUTO/\ldots} (qsqli);

\end{tikzpicture}
\caption{NFA para detección de SQL Injection (\texttt{pgsql-sqli-v2}).
  $q_0$=inicial; $q_!$=aceptación; abreviaturas: IS=INFORMATION\_SCHEMA,
  U\_A=UNION\_ALL, TAUTO=TAUTOLOGY, SEMI=SEMICOLON.}
\label{fig:nfa-sqli}
\end{figure}

%----------------------------------------------------------------------
\section{NFA — Escalada de Privilegios}
\noindent\textbf{Perfil:} \texttt{pgsql-privesc-v1}.

Este autómata reconoce la sentencia \texttt{ALTER USER/ROLE <nombre> [WITH] SUPERUSER},
usada para conceder acceso de superusuario en PostgreSQL. Tiene 6 estados y 7
transiciones, lo que lo hace altamente preciso con mínimos falsos positivos.

\begin{figure}[htbp]
\centering
\begin{tikzpicture}[
  ->, >=Stealth,
  node distance=2.6cm,
  every state/.style={draw, circle, minimum size=1.0cm, font=\small},
  accepting/.style={double},
  initial text=,
  auto
]

\node[state, initial]           (q0)     {$q_0$};
\node[state, right of=q0]       (qalt)   {$q_a$};
\node[state, right of=qalt]     (qtgt)   {$q_t$};
\node[state, right of=qtgt]     (qname)  {$q_n$};
\node[state, right of=qname]    (qwith)  {$q_w$};
\node[state, accepting,
      below of=qname, yshift=0.4cm]
                                 (qsuper) {\textbf{$q_s$}};

\draw (q0)    edge node[above]{\small ALTER}     (qalt);
\draw (qalt)  edge node[above]{\small USER}      (qtgt);
\draw (qalt)  edge[bend right=30] node[below]{\small ROLE} (qtgt);
\draw (qtgt)  edge node[above]{\small IDENT}     (qname);
\draw (qname) edge node[above]{\small WITH}      (qwith);
\draw (qname) edge node[right]{\small SUPERUSER} (qsuper);
\draw (qwith) edge node[right]{\small SUPERUSER} (qsuper);

\end{tikzpicture}
\caption{NFA para detección de Privilege Escalation (\texttt{pgsql-privesc-v1}).
  $q_0$=inicial; $q_s$=aceptación (Critical).}
\label{fig:nfa-privesc}
\end{figure}

%----------------------------------------------------------------------
\section{NFA — Exfiltración de Datos (\texttt{pgsql-exfil-v1})}

Detecta volcados masivos: \texttt{SELECT */cols FROM tabla} sin cláusula
\texttt{WHERE} ni \texttt{LIMIT}. La restricción negativa ($\mathit{absent} =
\{\tok{WHERE}, \tok{LIMIT}\}$) es el elemento clave.

\begin{figure}[htbp]
\centering
\begin{tikzpicture}[
  ->, >=Stealth,
  node distance=2.8cm,
  every state/.style={draw, circle, minimum size=1.0cm, font=\small},
  accepting/.style={double},
  initial text=,
  auto
]

\node[state, initial]          (q0)   {$q_0$};
\node[state, right of=q0]      (qsel) {$q_{sel}$};
\node[state, right of=qsel]    (qkf)  {$q_{kf}$};
\node[state, right of=qkf]     (qfr)  {$q_{fr}$};
\node[state, accepting,
      below of=qkf]            (qex)  {\textbf{$q_{ex}$}};

\draw (q0)   edge node[above]{\small SELECT}  (qsel);
\draw (qsel) edge[loop above] node{\small IDENT}   (qsel);
\draw (qsel) edge node[above]{\small STAR}    (qkf);
\draw (qsel) edge[bend right=30] node[below]{\small FROM}  (qfr);
\draw (qkf)  edge node[above]{\small FROM}    (qfr);
\draw (qfr)  edge node[right]{\small IDENT}   (qex);

\end{tikzpicture}
\caption{NFA para detección de Exfiltración (\texttt{pgsql-exfil-v1}).
  Aceptación requiere adicionalmente ausencia de \tok{WHERE} y \tok{LIMIT}.}
\label{fig:nfa-exfil}
\end{figure}

%======================================================================
\chapter{Tablas de Transiciones}
%======================================================================

\section{SQL Injection — \texttt{pgsql-sqli-v2}}

La siguiente tabla presenta la función de transición $\delta$ completa del
autómata de SQL Injection. Los estados se nombran abreviados para facilitar
la lectura.

\begin{center}
\small
\begin{longtable}{lll}
\caption{Tabla de transiciones — SQL Injection (\texttt{pgsql-sqli-v2})}
\label{tab:delta-sqli}\\
  \toprule
  \textbf{Estado origen} & \textbf{Símbolo} & \textbf{Estado destino} \\
  \midrule
  \endfirsthead
  \caption[]{Tabla de transiciones — SQL Injection (continuación)}\\
  \toprule
  \textbf{Estado origen} & \textbf{Símbolo} & \textbf{Estado destino} \\
  \midrule
  \endhead
  \bottomrule
  \endfoot
  \midrule
  $q_0$       & SELECT             & $q_{\text{prefix}}$ \\
  $q_0$       & TAUTOLOGY          & $q_{\text{sqli}}$   \\
  $q_0$       & SLEEP              & $q_{\text{sqli}}$   \\
  $q_0$       & INFORMATION\_SCHEMA& $q_{\text{sqli}}$   \\
  $q_0$       & OR                 & $q_{\text{or}}$     \\
  $q_0$       & UNION              & $q_{\text{union}}$  \\
  $q_0$       & UNION\_ALL         & $q_{\text{union}}$  \\
  \midrule
  $q_{\text{prefix}}$ & IDENT      & $q_{\text{prefix}}$ \\
  $q_{\text{prefix}}$ & STAR       & $q_{\text{prefix}}$ \\
  $q_{\text{prefix}}$ & FROM       & $q_{\text{tbl}}$    \\
  $q_{\text{prefix}}$ & SLEEP      & $q_{\text{sqli}}$   \\
  \midrule
  $q_{\text{tbl}}$ & IDENT         & $q_{\text{tbl}}$    \\
  $q_{\text{tbl}}$ & INFORMATION\_SCHEMA & $q_{\text{sqli}}$ \\
  $q_{\text{tbl}}$ & WHERE         & $q_{\text{where}}$  \\
  $q_{\text{tbl}}$ & SEMICOLON     & $q_{\text{sqli}}$   \\
  \midrule
  $q_{\text{where}}$ & IDENT       & $q_{\text{col}}$    \\
  $q_{\text{where}}$ & TAUTOLOGY   & $q_{\text{sqli}}$   \\
  $q_{\text{where}}$ & SEMICOLON   & $q_{\text{sqli}}$   \\
  \midrule
  $q_{\text{col}}$ & EQUALS        & $q_{\text{eq}}$     \\
  \midrule
  $q_{\text{eq}}$ & STRING         & $q_{\text{val}}$    \\
  $q_{\text{eq}}$ & NUMBER         & $q_{\text{val}}$    \\
  \midrule
  $q_{\text{val}}$ & OR            & $q_{\text{sqli}}$   \\
  $q_{\text{val}}$ & UNION         & $q_{\text{sqli}}$   \\
  $q_{\text{val}}$ & UNION\_ALL    & $q_{\text{sqli}}$   \\
  $q_{\text{val}}$ & TAUTOLOGY     & $q_{\text{sqli}}$   \\
  $q_{\text{val}}$ & SLEEP         & $q_{\text{sqli}}$   \\
  $q_{\text{val}}$ & INFORMATION\_SCHEMA & $q_{\text{sqli}}$ \\
  $q_{\text{val}}$ & SEMICOLON     & $q_{\text{sqli}}$   \\
  \midrule
  $q_{\text{or}}$ & TAUTOLOGY      & $q_{\text{sqli}}$   \\
  $q_{\text{or}}$ & NUMBER         & $q_{\text{sqli}}$   \\
  $q_{\text{or}}$ & STRING         & $q_{\text{sqli}}$   \\
  \midrule
  $q_{\text{union}}$ & SELECT      & $q_{\text{sqli}}$   \\
  $q_{\text{union}}$ & ALL         & $q_{\text{sqli}}$   \\
  \midrule
  $q_{\text{sqli}}$ & SELECT       & $q_{\text{sqli}}$   \\
  $q_{\text{sqli}}$ & TAUTOLOGY    & $q_{\text{sqli}}$   \\
  $q_{\text{sqli}}$ & UNION        & $q_{\text{sqli}}$   \\
  $q_{\text{sqli}}$ & UNION\_ALL   & $q_{\text{sqli}}$   \\
  $q_{\text{sqli}}$ & SLEEP        & $q_{\text{sqli}}$   \\
  \bottomrule
\end{longtable}
\end{center}

\section{Privilege Escalation — \texttt{pgsql-privesc-v1}}

\begin{table}[htbp]
\centering
\caption{Tabla de transiciones — Privilege Escalation (\texttt{pgsql-privesc-v1})}
\label{tab:delta-privesc}
\begin{tabular}{lll}
  \toprule
  \textbf{Estado origen} & \textbf{Símbolo} & \textbf{Estado destino} \\
  \midrule
  $q_0$          & ALTER     & $q_{\text{alter}}$  \\
  $q_{\text{alter}}$  & USER      & $q_{\text{target}}$ \\
  $q_{\text{alter}}$  & ROLE      & $q_{\text{target}}$ \\
  $q_{\text{target}}$ & IDENT     & $q_{\text{name}}$   \\
  $q_{\text{name}}$   & WITH      & $q_{\text{with}}$   \\
  $q_{\text{name}}$   & SUPERUSER & $q_{\text{super}}$  \\
  $q_{\text{with}}$   & SUPERUSER & $q_{\text{super}}$  \\
  \bottomrule
\end{tabular}
\end{table}

\section{Data Exfiltration — \texttt{pgsql-exfil-v1}}

\begin{table}[htbp]
\centering
\caption{Tabla de transiciones — Data Exfiltration (\texttt{pgsql-exfil-v1})}
\label{tab:delta-exfil}
\begin{tabular}{lll}
  \toprule
  \textbf{Estado origen} & \textbf{Símbolo} & \textbf{Estado destino} \\
  \midrule
  $q_0$               & SELECT & $q_{\text{sel}}$   \\
  $q_{\text{sel}}$    & IDENT  & $q_{\text{sel}}$   \\
  $q_{\text{sel}}$    & STAR   & $q_{\text{kw\_from}}$ \\
  $q_{\text{sel}}$    & FROM   & $q_{\text{from}}$  \\
  $q_{\text{kw\_from}}$ & FROM & $q_{\text{from}}$  \\
  $q_{\text{from}}$   & IDENT  & $q_{\text{ex}}$    \\
  \bottomrule
\end{tabular}
\end{table}

\begin{nota}
  Para el perfil de exfiltración, la tabla de transiciones define únicamente
  las transiciones positivas. La restricción negativa $\mathit{absent} =
  \{\texttt{WHERE}, \texttt{LIMIT}\}$ se evalúa \emph{post-facto}: si la cadena
  de tokens contiene alguno de estos dos tokens, la aceptación se cancela
  aunque el estado $q_{\text{ex}}$ sea alcanzado.
\end{nota}

%======================================================================
\chapter{Decisiones de Diseño Justificadas}
%======================================================================

\section{Uso de NFA en lugar de expresiones regulares simples}

\textbf{Decisión:} modelar cada patrón de ataque como un NFA serializado en
JSON en lugar de codificar expresiones regulares directamente en el código.

\textbf{Justificación:}
\begin{itemize}
  \item \emph{Extensibilidad sin recompilación:} agregar un nuevo perfil de
    amenaza requiere únicamente crear un archivo JSON y copiarlo al directorio
    \texttt{NFA/}. No se modifica ni compila el código C\#.
  \item \emph{Separación de responsabilidades:} el motor (\texttt{NfaEngine})
    es agnóstico al tipo de amenaza; los analistas de seguridad pueden definir
    o ajustar perfiles sin conocer el código fuente.
  \item \emph{Auditabilidad:} el estado del sistema de detección queda
    explicitamente documentado en archivos de texto (JSON) que pueden
    versionarse, revisarse y auditarse.
  \item \emph{Composabilidad:} múltiples perfiles se ejecutan en paralelo
    sobre la misma secuencia de tokens en $O(|w| \cdot |Q_{\text{total}}|)$.
\end{itemize}

\section{Tokenización en cuatro fases con marcador TAUTOLOGY}

\textbf{Decisión:} incorporar una fase de pre-análisis semántico (Fase~2) que
reemplaza expresiones tautológicas por el token especial \tok{TAUTOLOGY} antes
del escaneo léxico.

\textbf{Justificación:}
\begin{itemize}
  \item Un NFA que intentara reconocer directamente \texttt{1=1} o \texttt{'x'='x'}
    necesitaría transiciones para todos los posibles valores numéricos y cadenas,
    lo que resultaría en un autómata de tamaño exponencial.
  \item Al abstraer la tautología a un único token, el NFA permanece pequeño
    y lineal, y el reconocimiento es $O(|w|)$ en lugar de $O(|w| \cdot |\Sigma_{\text{val}}|)$.
  \item Las 5 expresiones regulares usadas para detectar tautologías están
    acotadas (longitudes máximas en cuantificadores) para prevenir ataques
    ReDoS (\emph{Regular expression Denial of Service}).
\end{itemize}

\section{Modo de reconocimiento de subcadena (\emph{substring match})}

\textbf{Decisión:} reinyectar $q_0$ en el conjunto activo en cada paso de
simulación, convirtiendo el NFA en un detector de subcadenas.

\textbf{Justificación:}
\begin{itemize}
  \item Las consultas SQL legítimas rara vez son secuencias puras del patrón
    de ataque; el payload malicioso suele estar embebido en una consulta
    más amplia (e.g., \texttt{SELECT id FROM t WHERE name='x' OR 1=1}).
  \item Exigir que el NFA reconozca la cadena \emph{completa} desde $q_0$
    generaría falsos negativos masivos.
  \item La reinyección de $q_0$ mantiene la complejidad lineal sin requerir
    la construcción explícita de un DFA o el uso de búsqueda por sufijos.
\end{itemize}

\section{Restricción negativa (\texttt{requireAbsentTokens})}

\textbf{Decisión:} añadir un mecanismo de post-condición negativa a los
perfiles NFA para invalidar la aceptación cuando ciertos tokens están presentes.

\textbf{Justificación:}
\begin{itemize}
  \item El perfil de exfiltración (\texttt{pgsql-exfil-v1}) detecta
    \texttt{SELECT * FROM tabla}. Sin la restricción negativa, cualquier consulta
    legítima como \texttt{SELECT * FROM config WHERE id=1} sería marcada como
    amenaza (falso positivo).
  \item Codificar la ausencia de \tok{WHERE} directamente en el NFA requeriría
    estados adicionales para cada token que \emph{no} sea WHERE, haciendo
    el autómata exponencialmente complejo.
  \item La verificación post-facto es $O(|\mathit{absent}|)$ y se realiza
    una sola vez por evaluación exitosa, sin impacto práctico en el rendimiento.
\end{itemize}

\section{Arquitectura de pipeline asíncrono}
\noindent\textbf{Componente:} \texttt{Channel\textless T\textgreater}.

\textbf{Decisión:} desacoplar la lectura de archivos, el parseo y la evaluación
NFA mediante un canal productor-consumidor (\texttt{System.Threading.Channels}).

\textbf{Justificación:}
\begin{itemize}
  \item PostgreSQL puede emitir rafagas de cientos de líneas por segundo bajo
    carga. Un modelo síncrono bloquearía la UI de WPF (que corre en el hilo
    principal) durante el procesamiento.
  \item \texttt{Channel\textless T\textgreater} proporciona backpressure natural: si el consumidor
    (NFA engine) va lento, el canal acumula sin perder entradas.
  \item El modelo \emph{async/await} de .NET 10 permite que la UI siga
    respondiendo mientras el motor procesa logs en segundo plano.
\end{itemize}

\section{Evaluación de perfiles en formato JSON con tipado fuerte}

\textbf{Decisión:} serializar los perfiles NFA en JSON con esquema fijo.
El tipo \texttt{AutomatonProfile} se deserializa con \texttt{System.Text.Json}.

\textbf{Justificación:}
\begin{itemize}
  \item JSON es un formato ubicuo, legible por humanos y soportado nativamente
    en .NET sin dependencias externas.
  \item El tipado fuerte (clases C\# con anotaciones \texttt{JsonPropertyName})
    detecta errores de configuración en tiempo de carga, no en tiempo de
    ejecución durante una detección.
  \item Los archivos JSON se copian al directorio de salida con la política
    \texttt{PreserveNewest}, garantizando que la versión del perfil coincida
    con la versión del binario.
\end{itemize}

%======================================================================
\chapter{Cadenas de Prueba}
%======================================================================

En esta sección se presentan cadenas de tokens de prueba con su trayectoria
de estados en el NFA correspondiente y el resultado esperado.

\section{Pruebas sobre SQL Injection (\texttt{pgsql-sqli-v2})}

\subsection{Caso P1 — Tautología directa}

\textbf{Consulta SQL:}
\begin{lstlisting}[style=json, numbers=none]
SELECT * FROM users WHERE id=1 OR 1=1
\end{lstlisting}

\textbf{Secuencia de tokens tras tokenización:}
\[
  \texttt{SELECT, STAR, FROM, IDENT, WHERE, IDENT, EQUALS, NUMBER, OR, TAUTOLOGY}
\]

\textbf{Traza de estados:}
\begin{center}
\begin{tabular}{llll}
  \toprule
  \textbf{Paso} & \textbf{Token} & \textbf{Estados activos} & \textbf{Acción} \\
  \midrule
  0  & —            & $\{q_0\}$                              & inicio \\
  1  & SELECT       & $\{q_0, q_{\text{prefix}}\}$           & $q_0 \xrightarrow{\text{SELECT}} q_{\text{prefix}}$ \\
  2  & STAR         & $\{q_0, q_{\text{prefix}}\}$           & loop \\
  3  & FROM         & $\{q_0, q_{\text{tbl}}\}$              & $q_{\text{prefix}} \xrightarrow{\text{FROM}} q_{\text{tbl}}$ \\
  4  & IDENT        & $\{q_0, q_{\text{tbl}}\}$              & loop \\
  5  & WHERE        & $\{q_0, q_{\text{where}}\}$            & $q_{\text{tbl}} \xrightarrow{\text{WHERE}} q_{\text{where}}$ \\
  6  & IDENT        & $\{q_0, q_{\text{col}}\}$              & $q_{\text{where}} \xrightarrow{\text{IDENT}} q_{\text{col}}$ \\
  7  & EQUALS       & $\{q_0, q_{\text{eq}}\}$               & $q_{\text{col}} \xrightarrow{\text{EQUALS}} q_{\text{eq}}$ \\
  8  & NUMBER       & $\{q_0, q_{\text{val}}\}$              & $q_{\text{eq}} \xrightarrow{\text{NUMBER}} q_{\text{val}}$ \\
  9  & OR           & $\{q_0, q_{\text{sqli}}, q_{\text{or}}\}$ & $q_{\text{val}} \xrightarrow{\text{OR}} q_{\text{sqli}}$ \\
  10 & TAUTOLOGY    & $\{q_0, q_{\text{sqli}}\}$             & $\mathbf{ACEPTA}$ \\
  \bottomrule
\end{tabular}
\end{center}

\textbf{Resultado:} \severity{danger}{AMENAZA} detectada. Tipo: \texttt{SQLI}, severidad: High.

\subsection{Caso P2 — UNION Injection}

\textbf{Consulta SQL:}
\begin{lstlisting}[style=json, numbers=none]
SELECT id FROM t UNION ALL SELECT password FROM pg_shadow
\end{lstlisting}

\textbf{Secuencia de tokens:}
\[
  \texttt{SELECT, IDENT, FROM, IDENT, UNION\_ALL, SELECT, IDENT, FROM, IDENT}
\]

\textbf{Traza abreviada:}
\begin{center}
\begin{tabular}{lll}
  \toprule
  \textbf{Token} & \textbf{Estados activos} & \textbf{Nota} \\
  \midrule
  SELECT    & $\{q_0, q_{\text{prefix}}\}$ & \\
  IDENT     & $\{q_0, q_{\text{prefix}}\}$ & loop \\
  FROM      & $\{q_0, q_{\text{tbl}}\}$    & \\
  IDENT     & $\{q_0, q_{\text{tbl}}\}$    & loop \\
  UNION\_ALL& $\{q_0, q_{\text{union}}\}$  & $q_0 \xrightarrow{\text{U\_ALL}} q_{\text{union}}$ \\
  SELECT    & $\{q_0, q_{\text{prefix}}, q_{\text{sqli}}\}$ & \textbf{ACEPTA} \\
  \bottomrule
\end{tabular}
\end{center}

\textbf{Resultado:} \severity{danger}{AMENAZA} — UNION Injection confirmado.

\subsection{Caso P3 — Consulta legítima (no debe activar alarma)}

\textbf{Consulta SQL:}
\begin{lstlisting}[style=json, numbers=none]
SELECT username, email FROM accounts WHERE user_id = 42
\end{lstlisting}

\textbf{Secuencia de tokens:}
\[
  \texttt{SELECT, IDENT, IDENT, FROM, IDENT, WHERE, IDENT, EQUALS, NUMBER}
\]

El autómata alcanza $q_{\text{val}}$ pero no recibe ningún token que
provoque transición a $q_{\text{sqli}}$. La cadena termina en
$\{q_0, q_{\text{val}}\}$ — $q_{\text{val}} \notin F$.

\textbf{Resultado:} \severity{success}{BENIGNA} — sin amenaza detectada. \checkmark

\section{Pruebas sobre Privilege Escalation}
\noindent\textbf{Perfil:} \texttt{pgsql-privesc-v1}.

\subsection{Caso P4 — Escalada directa}

\textbf{Consulta:} \texttt{ALTER USER administrador SUPERUSER}

\textbf{Tokens:} \texttt{ALTER, USER, IDENT, SUPERUSER}

\begin{center}
\begin{tabular}{lll}
  \toprule
  \textbf{Token} & \textbf{Estados activos} & \\
  \midrule
  ALTER     & $\{q_0, q_{\text{alter}}\}$  & \\
  USER      & $\{q_0, q_{\text{target}}\}$ & \\
  IDENT     & $\{q_0, q_{\text{name}}\}$   & \\
  SUPERUSER & $\{q_0, q_{\text{super}}\}$  & \textbf{ACEPTA} \\
  \bottomrule
\end{tabular}
\end{center}

\textbf{Resultado:} \severity{danger}{AMENAZA CRÍTICA} — Privilege Escalation.

\subsection{Caso P5 — Con cláusula WITH}

\textbf{Consulta:} \texttt{ALTER ROLE devuser WITH SUPERUSER}

\textbf{Tokens:} \texttt{ALTER, ROLE, IDENT, WITH, SUPERUSER}

Traza: $q_0 \to q_{\text{alter}} \to q_{\text{target}} \to q_{\text{name}}
\to q_{\text{with}} \to q_{\text{super}}$.

\textbf{Resultado:} \severity{danger}{AMENAZA CRÍTICA} detectada también por
la ruta alternativa \texttt{WITH}.

\section{Pruebas sobre Data Exfiltration}
\noindent\textbf{Perfil:} \texttt{pgsql-exfil-v1}.

\subsection{Caso P6 — Volcado masivo}

\textbf{Consulta:} \texttt{SELECT * FROM clientes}

\textbf{Tokens:} \texttt{SELECT, STAR, FROM, IDENT}

Ningún token en $\mathit{absent} = \{\texttt{WHERE, LIMIT}\}$ presente.
Autómata alcanza $q_{\text{ex}}$.

\textbf{Resultado:} \severity{danger}{AMENAZA} — Data Exfiltration.

\subsection{Caso P7 — SELECT * legítimo con WHERE}

\textbf{Consulta:} \texttt{SELECT * FROM pedidos WHERE estado='activo' LIMIT 100}

Tokens incluyen \tok{WHERE} y \tok{LIMIT}. El autómata alcanza $q_{\text{ex}}$
pero la restricción negativa invalida la aceptación.

\textbf{Resultado:} \severity{success}{BENIGNA} — restricción negativa activa. \checkmark

\section{Resumen de Cadenas de Prueba}

\begin{table}[htbp]
\centering
\caption{Resumen de resultados de pruebas}
\label{tab:pruebas}
\begin{tabularx}{\textwidth}{clXcc}
  \toprule
  \textbf{ID} & \textbf{Perfil} & \textbf{Patrón} & \textbf{Esperado} & \textbf{Obtenido} \\
  \midrule
  P1 & SQLI      & Tautología OR 1=1            & AMENAZA  & \color{success}\checkmark \\
  P2 & SQLI      & UNION ALL SELECT pg\_shadow  & AMENAZA  & \color{success}\checkmark \\
  P3 & SQLI      & Consulta legítima con WHERE  & BENIGNA  & \color{success}\checkmark \\
  P4 & PRIVESC   & ALTER USER SUPERUSER directo & AMENAZA  & \color{success}\checkmark \\
  P5 & PRIVESC   & ALTER ROLE WITH SUPERUSER    & AMENAZA  & \color{success}\checkmark \\
  P6 & EXFIL     & SELECT * FROM sin WHERE      & AMENAZA  & \color{success}\checkmark \\
  P7 & EXFIL     & SELECT * con WHERE y LIMIT   & BENIGNA  & \color{success}\checkmark \\
  \bottomrule
\end{tabularx}
\end{table}

Todos los casos producen el resultado esperado, validando la corrección de los
autómatas y la eficacia de los mecanismos auxiliares (tautología, restricción
negativa, modo substring).

%======================================================================
\chapter{Conclusiones}
%======================================================================

\section{Síntesis del sistema}

LogGuard V2 implementa un sistema de detección de intrusiones en tiempo real
para PostgreSQL basado en la teoría formal de autómatas finitos. La arquitectura
puede resumirse en tres pilares:

\begin{enumerate}
  \item \textbf{Análisis léxico en dos niveles:} el parseo de logs de PostgreSQL
    produce entradas estructuradas; la tokenización SQL de cuatro fases transforma
    sentencias SQL en secuencias sobre un alfabeto controlado de más de 70 tokens.
  \item \textbf{Detección mediante NFA:} seis perfiles NFA serializados en JSON
    reconocen subcadenas de tokens que corresponden a patrones de ataque
    conocidos (SQL Injection, Exfiltración, Escalada de Privilegios, Fuerza
    Bruta, Descubrimiento y Inyección Temporal).
  \item \textbf{Arquitectura reactiva:} un pipeline asíncrono basado en
    \texttt{Channel\textless T\textgreater} permite monitorear logs con baja latencia sin bloquear
    la interfaz gráfica.
\end{enumerate}

\section{Fortalezas del diseño}

\begin{itemize}
  \item \emph{Extensibilidad:} nuevos vectores de ataque se incorporan añadiendo
    archivos JSON sin modificar el código del motor.
  \item \emph{Resistencia a evasión:} la normalización (Fase~1) y la detección
    semántica de tautologías (Fase~2) contrarrestan técnicas comunes de
    ofuscación SQL.
  \item \emph{Rendimiento predecible:} la complejidad lineal $O(|w| \cdot |Q|)$
    del simulador de subconjuntos garantiza tiempos de respuesta acotados
    independientemente del contenido de la consulta.
  \item \emph{Precisión ajustable:} la restricción negativa reduce falsos
    positivos sin aumentar la complejidad del autómata.
\end{itemize}

\section{Limitaciones y trabajo futuro}

\begin{itemize}
  \item \emph{Cobertura de vocabulario:} el sistema reconoce dialecto PostgreSQL.
    Extensión a otros dialectos (MySQL, MSSQL, Oracle) requeriría perfiles NFA
    adicionales y ajustes al tokenizador.
  \item \emph{Ataques de contexto múltiple:} ataques que distribuyen el payload
    en múltiples consultas (e.g., inyección por batches) no son detectados por
    los perfiles actuales, que operan consulta por consulta.
  \item \emph{Aprendizaje automático:} los perfiles NFA son diseñados manualmente.
    Un sistema de inferencia automática de autómatas (e.g., L*) a partir de
    tráfico etiquetado aumentaría la cobertura de detección.
  \item \emph{Explotación de SLEEP:} la detección de inyección basada en tiempo
    mediante el token \tok{SLEEP} no considera el tiempo real de ejecución
    registrado en las líneas \texttt{duration:}, que sería un indicador más
    preciso.
\end{itemize}

\section{Conclusión final}

LogGuard V2 demuestra que la teoría formal de autómatas tiene aplicación
directa y práctica en sistemas de seguridad modernos. La elección de NFA como
motor de detección ofrece un balance óptimo entre expresividad formal,
rendimiento computacional y mantenibilidad operacional. El sistema es capaz de
detectar en tiempo real las principales clases de ataques SQL sobre PostgreSQL,
con tasas de falsos positivos bajas gracias a los mecanismos de restricción
negativa y normalización semántica.

\bigskip
\begin{center}
\rule{8cm}{0.4pt}\\[0.3cm]
{\small Adolfo Toledo — LogGuard V2 — Junio 2026}
\end{center}

%======================================================================
\end{document}
%======================================================================
